\documentclass[10pt, aspectratio=169]{beamer}
\input{../../_en_preambulo_beamer}
% Comandos y macros estadísticas compartidas para las presentaciones Beamer en inglés (Metropolis)

% Conjuntos numéricos básicos
\newcommand{\R}{\mathbb{R}}
\newcommand{\N}{\mathbb{N}}
\newcommand{\Q}{\mathbb{Q}}
\newcommand{\Z}{\mathbb{Z}}
\newcommand{\set}[1]{\left\{ #1 \right\}}
\newcommand{\abs}[1]{\left|#1\right|}
\newcommand{\norm}[1]{\left\| #1 \right\|}

% Comandos estadísticos y probabilísticos del curso
\newcommand{\Var}{\operatorname{Var}}
\newcommand{\cov}{\operatorname{Cov}}
\newcommand{\corr}{\rho}
\newcommand{\comb}[2]{\begin{pmatrix} #1 \\ #2 \end{pmatrix}}
\newcommand{\s}{\sigma}
\newcommand{\particion}{\mathfrak{P}}
\newcommand{\card}[1]{\#\left( #1 \right)}
\newcommand{\seq}[3]{\set{#1}_{#2}^{#3}}
\newcommand{\E}{\mathbb{E}}
\newcommand{\Prob}{\mathbb{P}}

% Entornos pedagógicos y cajas en inglés académico natural (soportando tanto nombres en inglés como en español)
\newenvironment{discussion}{%
  \begin{alertblock}{Discussion Question}%
}{%
  \end{alertblock}%
}
\newenvironment{discusion}{%
  \begin{alertblock}{Discussion Question}%
}{%
  \end{alertblock}%
}

\newenvironment{reflection}{%
  \begin{alertblock}{Classroom Reflection}%
}{%
  \end{alertblock}%
}
\newenvironment{reflexion}{%
  \begin{alertblock}{Classroom Reflection}%
}{%
  \end{alertblock}%
}

\newenvironment{paradox}{%
  \begin{alertblock}{Exploratory Paradox}%
}{%
  \end{alertblock}%
}
\newenvironment{paradoja}{%
  \begin{alertblock}{Exploratory Paradox}%
}{%
  \end{alertblock}%
}

\newenvironment{motivation}{%
  \begin{exampleblock}{Motivation: Data Science in Practice}%
}{%
  \end{exampleblock}%
}
\newenvironment{motivacion}{%
  \begin{exampleblock}{Motivation: Data Science in Practice}%
}{%
  \end{exampleblock}%
}

\newenvironment{quickchallenge}{%
  \begin{block}{Quick Team Challenge}%
}{%
  \end{block}%
}
\newenvironment{retoclase}{%
  \begin{block}{Quick Team Challenge}%
}{%
  \end{block}%
}


\title[ANOVA Assumptions]{Section 08.02: ANOVA Assumptions and Diagnostics}
\subtitle{Homoscedasticity, Normality, Independence, and Nonparametric Alternatives}
\author[J. Castillo Colmenares]{Juliho Castillo Colmenares}
\institute{Tecnológico de Monterrey}
\date{\vspace{-1.5cm}}

\begin{document}

% Slide 1: Title
\begin{frame}[plain]
  \titlepage
\end{frame}

% Slide 2: Roadmap
\begin{frame}{Roadmap --- Chapter 08: Elements of Experimental Design (ANOVA)}
  \footnotesize
  Where are we in the modular development of this chapter? \pause
  \begin{itemize}\itemsep=0.08em
    \item \textbf{08.01} One-Way Analysis of Variance (ANOVA) and the $F$-Test \pause
    \item \textbf{08.02} ANOVA Assumptions, the Levene/Bartlett Tests, and Diagnostics \emph{(Today --- Chapter Closing)}
  \end{itemize}
  \vspace{-0.05cm}
  \begin{block}{Session Objective}
    Rigorously audit the three assumptions underlying the validity of all the machinery built in Section 08.01 --- homoscedasticity, normality, and independence of the errors --- and present the methodological alternatives (Welch's ANOVA, Kruskal-Wallis) when those assumptions fail.
  \end{block}
\end{frame}

% Slide 3: Motivation
\begin{frame}{Why Audit the Assumptions?}
  \footnotesize
  \begin{motivacion}
    The statistic $F=\text{CMTR}/\text{CME}$ only follows an \emph{exact} $F_{k-1,N-k}$ distribution if the underlying parametric model ($\epsilon_{ij}\sim N(0,\sigma^2)$ i.i.d.) is correct. An ANOVA computed on data that violate these assumptions can produce completely wrong decisions, even with flawless arithmetic.
  \end{motivacion}

  \vspace{0.15cm}
  \pause
  \begin{itemize}\itemsep=0.1cm
    \item Every data scientist must audit these assumptions \textbf{before} reporting ANOVA results in production or research. \pause
    \item The diagnostic tests mostly reuse the same variance-decomposition machinery from Section 08.01.
  \end{itemize}
\end{frame}

% Slide 4: The three assumptions
\begin{frame}{The Three Pillars of the Parametric Model}
  \footnotesize
  \begin{itemize}\itemsep=0.12cm
    \item \textbf{Homoscedasticity:} $\sigma_1^2=\sigma_2^2=\ldots=\sigma_k^2$. ANOVA is robust to moderate violations \emph{only if} sample sizes are equal. \pause
    \item \textbf{Normality of errors:} $\epsilon_{ij}\sim N(0,\sigma^2)$. Audited via Q-Q plots and formal tests on the residuals. \pause
    \item \textbf{Independence:} $\text{Cov}(\epsilon_{ij},\epsilon_{lm})=0$. The most critical assumption: its violation completely invalidates the $F$-test and cannot be fixed with transformations; it is guaranteed \emph{a priori} through randomization.
  \end{itemize}
\end{frame}

% Slide 5: Homoscedasticity
\begin{frame}{Verifying Homoscedasticity: Bartlett and Levene}
  \footnotesize
  \begin{block}{Bartlett's Test}
    Sensitive to normality. Tests $H_0:\sigma_1^2=\ldots=\sigma_k^2$ via a statistic based on the weighted ratio of sample variances, $\sim\chi^2_{k-1}$.
  \end{block}
  \pause

  \vspace{0.1cm}
  \begin{alertblock}{Levene's Test (mean-centered) and Brown-Forsythe (median)}
    Modern gold standard: applies the full one-way ANOVA from Section 08.01 to the absolute deviations $Z_{ij}=|Y_{ij}-\bar Y_{i.}|$. Robust to non-normality.
  \end{alertblock}
\end{frame}

% Slide 6: Normality
\begin{frame}{Diagnosing Residual Normality}
  \footnotesize
  \begin{block}{Qualitative Audit}
    The normal probability plot (\textbf{Q-Q Plot}) of the standardized residuals $\hat\epsilon_{ij}=Y_{ij}-\bar Y_{i.}$ visually reveals systematic departures from normality.
  \end{block}
  \pause

  \vspace{0.1cm}
  \begin{alertblock}{Quantitative Audit}
    The \textbf{Shapiro-Wilk Test} (or Kolmogorov-Smirnov) formally tests $H_0:$ the residuals come from a normal population. It is the reference test for small to moderate samples.
  \end{alertblock}
\end{frame}

% Slide 7: Alternatives
\begin{frame}{When Assumptions Fail: Methodological Alternatives}
  \footnotesize
  \begin{itemize}\itemsep=0.1cm
    \item \textbf{Homoscedasticity violated:} Welch's ANOVA (adjusts denominator degrees of freedom by weighting $S_i^2/n_i$) or Box-Cox transformations ($\log Y$, $\sqrt{Y}$). \pause
    \item \textbf{Normality violated (extreme skew, outliers):} \textbf{Kruskal-Wallis Test}, the nonparametric alternative based on ranks instead of means. \pause
    \item \textbf{Independence violated:} no algebraic after-the-fact correction exists; requires time series or longitudinal mixed-effects models.
  \end{itemize}
\end{frame}

% Slide 8: Python Lab (Block 1)
\begin{frame}[fragile]{Python Lab: Bartlett and Levene}
  \scriptsize
  Homoscedasticity audit for 3 CI/CD pipelines, verified against \texttt{scipy.stats.bartlett} and \texttt{scipy.stats.levene} (\texttt{08.02\_anova\_assumptions.py}):
  {\lstset{basicstyle=\fontsize{3.5pt}{4.2pt}\selectfont\ttfamily,aboveskip=0.05em,belowskip=0.05em}
  \lstinputlisting[firstline=17, lastline=38]{../../code/08_diseno_experimentos/08.02_anova_assumptions.py}}
\end{frame}

% Slide 9: Python Lab (Block 2)
\begin{frame}[fragile]{Python Lab: Shapiro-Wilk on Residuals}
  \scriptsize
  Fitting a 3-group ANOVA and auditing the normality of its residuals:
  {\lstset{basicstyle=\fontsize{3.8pt}{4.5pt}\selectfont\ttfamily,aboveskip=0.05em,belowskip=0.05em}
  \lstinputlisting[firstline=41, lastline=52]{../../code/08_diseno_experimentos/08.02_anova_assumptions.py}}
\end{frame}

% Slide 10: Python Lab (Block 3)
\begin{frame}[fragile]{Python Lab: Kruskal-Wallis}
  \scriptsize
  Comparing the parametric ANOVA against Kruskal-Wallis on skewed (exponential) data:
  {\lstset{basicstyle=\fontsize{3.6pt}{4.3pt}\selectfont\ttfamily,aboveskip=0.05em,belowskip=0.05em}
  \lstinputlisting[firstline=55, lastline=70]{../../code/08_diseno_experimentos/08.02_anova_assumptions.py}}
\end{frame}

% Slide 11: Terminal Output
\begin{frame}[fragile]{Python Lab: Terminal Output}
  \scriptsize
  Standard output produced when running the lab script:
  \vspace{0.02cm}
  \begin{block}{Standard Console Output}
    \fontsize{3.8pt}{4.6pt}\selectfont
    \begin{verbatim}
=== Block 1: Bartlett and Levene ===
Bartlett: B=4.4660, p=0.1072
Levene (manual): W=2.9880, p=0.0885 (matches scipy.stats.levene center='mean')

=== Block 2: Shapiro-Wilk ===
W=0.9781, p=0.7721 -> no evidence against normality

=== Block 3: Kruskal-Wallis ===
Parametric ANOVA: F=3.5514, p=0.0376
Kruskal-Wallis:    H=11.9985, p=0.0025
    \end{verbatim}
  \end{block}
\end{frame}

% Slide 12A: Exercise Level 1 (Statement)
\begin{frame}{In-Class Exercise --- Level 1: Fundamental (1/2)}
  \footnotesize
  \begin{block}{Problem (Statement, Problem 8.6.11)}
    An analyst audits the homoscedasticity of $k=4$ configurations of a distributed pipeline, $n_i=8$ each, with $S_1^2=4.0$, $S_2^2=4.5$, $S_3^2=5.2$, $S_4^2=4.8$. Compute the Bartlett statistic and compare it against $\chi^2_{0.05,3}=7.815$.
  \end{block}

  \vspace{0.15cm}
  \pause
  \begin{alertblock}{Setup}
    First compute the pooled variance $S_p^2 = \sum(n_i-1)S_i^2/(N-k)$.
  \end{alertblock}
\end{frame}

% Slide 12B: Exercise Level 1 (Resolution)
\begin{frame}{In-Class Exercise --- Level 1: Fundamental (2/2)}
  \scriptsize
  We compute the pooled variance and the Bartlett statistic: \pause

  \vspace{0.05cm}
  \begin{block}{Calculation}
    $$S_p^2=\frac{7(18.5)}{28}=4.625, \qquad B \approx \frac{28\ln(4.625)-42.754}{1.0595} \approx 0.121$$
  \end{block}

  \vspace{0.05cm}
  \pause
  \begin{alertblock}{Interpretation}
    Since $0.121 \ll 7.815$, \textbf{we do not reject $H_0$}: the four variances are statistically compatible with homoscedasticity.
  \end{alertblock}
\end{frame}

% Slide 13A: Exercise Level 2 (Statement)
\begin{frame}{In-Class Exercise --- Level 2: Operational (1/2)}
  \footnotesize
  \begin{block}{Problem (Statement, Problem 8.6.12)}
    Three CI/CD pipelines with $n=5$ runs: A$=\{40,42,38,44,41\}$, B$=\{50,55,48,52,60\}$, C$=\{41,40,42,39,43\}$. Compute $Z_{ij}=|Y_{ij}-\bar Y_{i.}|$ and apply the Section 08.01 ANOVA to the $Z_{ij}$ (Levene) at $\alpha=0.05$.
  \end{block}

  \vspace{0.15cm}
  \pause
  \begin{alertblock}{Strategy}
    \scriptsize
    Means: $\bar Y_{A.}=41$, $\bar Y_{B.}=53$, $\bar Y_{C.}=41$; compare $W$ against $F_{0.05,2,12}=3.885$.
  \end{alertblock}
\end{frame}

% Slide 13B: Exercise Level 2 (Resolution)
\begin{frame}{In-Class Exercise --- Level 2: Operational (2/2)}
  \scriptsize
  We apply the sum-of-squares decomposition to the $Z_{ij}$: \pause

  \vspace{0.05cm}
  \begin{block}{Calculation}
    $$\text{SCTR}_Z\approx16.533,\ \text{SCE}_Z=33.2, \qquad \text{CMTR}_Z\approx8.267,\ \text{CME}_Z\approx2.767, \qquad W\approx2.988$$
  \end{block}

  \vspace{0.05cm}
  \pause
  \begin{alertblock}{Interpretation}
    \scriptsize
    Since $2.988<3.885$, \textbf{we do not reject $H_0$}: although Pipeline B looks more dispersed visually, $n=5$ per group is insufficient to declare significant heteroscedasticity.
  \end{alertblock}
\end{frame}

% Slide 14A: Exercise Level 3 (Statement)
\begin{frame}{In-Class Exercise --- Level 3: Analytical (1/2)}
  \footnotesize
  \begin{block}{Problem (Statement, Problem 8.6.7)}
    An ANOVA with $k=3$ cache configurations, $n_i=6$ each, has $S_1^2=1.2$, $S_2^2=1.5$, $S_3^2=8.9$. An analyst claims ``ANOVA is robust as long as sample sizes are equal.'' Is this correct? Which procedure and alternative apply?
  \end{block}

  \vspace{0.15cm}
  \pause
  \begin{alertblock}{Strategy}
    \scriptsize
    Compute the ratio of the extreme variances and formally contrast it via Levene.
  \end{alertblock}
\end{frame}

% Slide 14B: Exercise Level 3 (Resolution)
\begin{frame}{In-Class Exercise --- Level 3: Analytical (2/2)}
  \scriptsize
  We evaluate the magnitude of the variance disparity: \pause

  \vspace{0.05cm}
  \begin{block}{Calculation}
    $$S_3^2/S_1^2 = 8.9/1.2 \approx 7.41$$
  \end{block}

  \vspace{0.05cm}
  \pause
  \begin{alertblock}{Interpretation}
    \scriptsize
    The claim is \textbf{incomplete}: ANOVA's robustness does not tolerate extreme disparities like $7.4:1$. Levene's Test will decisively reject $H_0$; the analyst should use \textbf{Welch's ANOVA} or a log transformation.
  \end{alertblock}
\end{frame}

% Slide 15A: Exercise Level 4 (Statement)
\begin{frame}{In-Class Exercise --- Level 4: Challenging (1/2)}
  \footnotesize
  \begin{block}{Problem --- Levene as ANOVA (Statement, Problem 8.6.13)}
    Prove that the mean-centered Levene Test is identical to applying the Section 08.01 one-way ANOVA to $Z_{ij}=|Y_{ij}-\bar Y_{i.}|$: show that $W\equiv\text{CMTR}_Z/\text{CME}_Z$.
  \end{block}

  \vspace{0.1cm}
  \pause
  \begin{alertblock}{Strategy}
    \scriptsize
    Apply the Section 08.01 sum-of-squares partition theorem directly to $Z$ instead of $Y$.
  \end{alertblock}
\end{frame}

% Slide 15B: Exercise Level 4 (Resolution)
\begin{frame}{In-Class Exercise --- Level 4: Challenging (2/2)}
  \scriptsize
  We reuse the entire ANOVA machinery on the transformed variable: \pause

  \vspace{0.05cm}
  \begin{block}{Calculation}
    \scriptsize
    $$\text{SCT}_Z=\text{SCTR}_Z+\text{SCE}_Z \implies F_Z=\frac{\text{SCTR}_Z/(k-1)}{\text{SCE}_Z/(N-k)}=\frac{N-k}{k-1}\cdot\frac{\sum_i n_i(\bar Z_{i.}-\bar Z_{..})^2}{\sum_i\sum_j(Z_{ij}-\bar Z_{i.})^2}\equiv W$$
  \end{block}

  \vspace{0.05cm}
  \pause
  \begin{alertblock}{Interpretation}
    \scriptsize
    $W$ and $F_Z$ are algebraically identical: Levene is not a new test but the same ANOVA applied to a re-expressed variable. Auditing homoscedasticity requires no tools beyond those already built in 08.01.
  \end{alertblock}
\end{frame}

% Slide 16: Closing
\begin{frame}{Section and Chapter Closing: ANOVA Assumptions}
  \begin{reflexion}
    With this section we close Chapter 08: we built one-way ANOVA as a generalization of mean comparison, and now rigorously audit the assumptions that guarantee its validity. No ANOVA result is trustworthy without verifying homoscedasticity, normality, and independence of the errors.
  \end{reflexion}

  \vspace{0.2cm}
  \pause
  \begin{block}{Key Takeaways}
    \begin{itemize}\itemsep=0.08cm
      \item \textbf{Levene $\equiv$ ANOVA on deviations:} no new machinery, just a variable transformation.
      \item \textbf{Conditional robustness:} ANOVA tolerates moderate heteroscedasticity only with equal sample sizes.
      \item \textbf{Alternatives:} Welch for heteroscedasticity, Kruskal-Wallis for severe non-normality.
    \end{itemize}
  \end{block}
\end{frame}

% Slide 17: Modular Perspective
\begin{frame}{Modular Perspective --- The Next Challenge}
  \scriptsize
  \begin{retoclase}
    With Chapter 08 complete, we have mastered comparing $k$ means through a single categorical factor. Chapter 09 (Simple and Multiple Linear Regression) generalizes this logic to \textbf{continuous} predictor variables: instead of comparing discrete groups, we will model the functional relationship between a response variable and one or more regressors via Ordinary Least Squares (OLS).
  \end{retoclase}

  \vspace{0.08cm}
  \pause
  \begin{alertblock}{Recommended Preparation}
    \begin{itemize}\itemsep=0.06cm
      \item Review the geometric interpretation of the sum-of-squares decomposition $\text{SCT}=\text{SCTR}+\text{SCE}$: in regression it will reappear as $\text{SCT}=\text{SCReg}+\text{SCE}$.
      \item Reflect on how the ANOVA coefficient of determination $R^2$ (Section 08.01) anticipates the $R^2$ of a regression model.
    \end{itemize}
  \end{alertblock}
\end{frame}

\end{document}
